\documentclass[a4paper, 12pt]{article}

% \usepackage[portuges]{babel}
\usepackage[utf8]{inputenc}
\usepackage{subfig}
\usepackage[colorinlistoftodos]{todonotes}

\usepackage[english]{babel}
\usepackage{indentfirst,setspace,subcaption}
\usepackage{amsmath,amssymb,graphicx,xcolor,url}
\usepackage{fancyhdr,tocbasic,titlesec,minted,listings}
\usepackage[a4paper,margin=24mm]{geometry}
\usepackage[skip=10pt plus1pt, indent=20pt]{parskip}
\usepackage[colorlinks=true,allcolors=blue,urlcolor=magenta]{hyperref}

\usepackage{indentfirst}
\usepackage{verbatim}
\usepackage{textcomp}
\usepackage{gensymb}

\usepackage{relsize}

\usepackage{lipsum}% http://ctan.org/pkg/lipsum
\usepackage{xcolor}% http://ctan.org/pkg/xcolor
\usepackage{xparse}% http://ctan.org/pkg/xparse
\NewDocumentCommand{\myrule}{O{1pt} O{2pt} O{black}}{%
  \par\nobreak% don't break a page here
  \kern\the\prevdepth% don't take into account the depth of the preceding line
  \kern#2 % space before the rule
  {\color{#3}\hrule height #1 width\hsize} % the rule
  \kern#2 % space after the rule
  \nointerlineskip% no additional space after the rule
}
\usepackage[section]{placeins}

\usepackage{booktabs}
\usepackage{colortbl}%
   \newcommand{\myrowcolour}{\rowcolor[gray]{0.925}}
   

\pagestyle{fancy}
\setlength{\headheight}{15pt}
\fancyhf{}
\fancyhead[R]{\nouppercase\rightmark\hfill~Exercise I Report}
\fancyfoot[C]{\hfill\thepage\hfill}

\DeclareTOCStyleEntry[
  indent=12pt,
  level=1
]{largetocline}{section}

\setlength{\marginparwidth}{2cm} 
\begin{document}

\begin{titlepage}
	\begin{center}
		\textbf{\LARGE Vietnam National University,}\\[0.5cm]
		\textbf{\LARGE Ho Chi Minh City}\\[0.5cm]
		\vspace{20pt}
		\textbf{\large UNIVERSITY OF SCIENCE}\\[0.2cm]
		\textbf{\large FACULTY OF INFORMATION TECHNOLOGY}\\[0.2cm]
		\vspace{20pt}
		\includegraphics[width=0.5\textwidth,keepaspectratio]{./images/logo.png}

		\par
		\vspace{20pt}
		\textbf{\Large [CSC10004] Data Structure And Algorithms}\\
		\vspace{15pt}
		\myrule[1pt][7pt]
		\textbf{\LARGE EXERCISE 2 REPORT}\\
		\vspace{15pt}
		\textbf{\large Implementing Hash Table from scratch}\\
		\vspace{10pt}
		\myrule[1pt][7pt]
		\vspace{25pt}
		\textbf{\large Student Name \hspace{20pt} Student ID}\\
		Bui Minh Duy \hspace{45pt} 23127040 \\

		\vspace{45pt}
		\textbf {\large Lecturer In Charge}\\[0.2cm]
		\Large {Bui Duy Dang}\\[0.1cm]
		\Large {Nguyen Thanh Tinh}\\[0.1cm]
	\end{center}

	\par
	\vfill
	\begin{center}
		\text{17/08/2024}\\
	\end{center}

\end{titlepage}

\tableofcontents\thispagestyle{empty}

\pagebreak
\section{Experiment Details}


\pagebreak
\section{Recursive versions}


\subsection{Conclusion}
\begin{itemize}
	\item Though Loop version and Recursive Version have the same theoretical time complexity of O(n), the actual execution time shows that Loop version is faster in all cases.
	\item When working with large data sets, the Loop version still executes as normal, but the Recursive version might encounter Stack Overflow.
\end{itemize}

\pagebreak
\section{Self-evaluation}
\begin{center}
	\renewcommand{\arraystretch}{1.5}
	\begin{tabular}{|l|p{\dimexpr0.6\linewidth-2\tabcolsep}|c|}
		\hline
		\textbf{No.} & \textbf{Details}   & \textbf{Score} \\ \hline
		1            & Stack Array        & 100\%          \\ \hline
		2            & Stack Linked List  & 100\%          \\ \hline
		3            & Queue Array        & 100\%          \\ \hline
		4            & Queue Linked List  & 100\%          \\ \hline
		5            & Recursive versions & 100\%          \\ \hline
		6            & Report             & 100\%          \\ \hline
		             & \textbf{Overall}   & \textbf{100\%} \\ \hline
	\end{tabular}
\end{center}

\section{Exercise Feedback}
\subsection{What I learned}
\begin{flushleft}
	\begin{itemize}
		\item Better understanding about the internal processing of stack and queue.
		\item How to use function templates in C++.
	\end{itemize}
\end{flushleft}
\subsection{What was difficult}
\begin{flushleft}
	\begin{itemize}
		\item The approach of using recursion is quite challenging, it took a significant amount of time and effort to complete the set-up.
	\end{itemize}
\end{flushleft}

\end{document}
